\chapter{22}



«Ey, Leute, jetzt müssen wir anfangen, sonst wird’s zu spät,» rief Meo
in die Runde, «sonst kommt auch schon Anton rüber und beschwert sich.»

Dazu verdrehte Meo die Augen und seine Worte kamen für Louis einer
Erlösung gleich. Meo schnitt ein grimmiges Gesicht, hielt die Hand hoch
und brabbelte, mit verstellter Stimme.

«Es ist längst Mitternacht durch, Ruhe jetzt, nicht zum Aushalten, diese
laute Musik.»

Sie lachten hell auf. Meos Mutter legte demonstrativ ihren Zeigefinger
vor ihre Lippen.

«Pass auf, Meo, sonst hört der dich,» sie klang beschwichtigend und
schmunzelte, «so ist der halt, der Anton,» sie verwarf die Hände, «den
können wir nicht ändern.»

\sectionbreak

Sie waren bereits alle aufgestanden. Louis’ Herz klopfte wild. Diesmal
angenehm. Er blickte erwartungsvoll um sich. Vanessa, Meo und seine
Eltern begannen zielsicher zwischen Haus und Garten hin und her zu
wirbeln, als folgten sie einem Plan. Das Schauspiel erinnerte Louis an
den Umbau einer Bühne zwischen zwei Theaterszenen. Er schaffte es kaum,
den raschen Bewegungen hinterher zu schauen.

Mit dem anhaltend hohen Tempo wurde das Geschirr abgetragen, stellten
sie die Gartenmöbel um und arrangierten Stühle in einem losen Kreis
inmitten des kleinen Platzes. Meo hatte einen beeindruckenden Verstärker
angeschleppt. Sein Vater legte ein Saxofon auf einen der Stühle. Dann
begann er mit Hingabe Kerzen auf hohen, schlanken Holzklötzen
anzuzünden. Vanessa packte die Gitarren aus. Louis stand perplex
mittendrin und gelegentlich im Weg.

Er war fasziniert, wie eingespielt und zielsicher die Vier einen
behaglichen Ort schafften.

«Du kannst dich in Ruhe vorbereiten, Ciro, wir machen hier alles klar.
Du kannst es bestimmt nicht erwarten, deine Gitarre auszupacken,» rief
Meo Louis aufgeregt zu, während er die Kabelkiste durchwühlte.

Er hatte ihn Ciro genannt. Bereits mehrmals. Es funktionierte.

Louis nickte. Er drehte sich um. Eine gute Gelegenheit, sich Adresse und
Telefonnummer des Artikels auf dem Küchentisch abzuschreiben. Niemand
bemerkte ihn, als er danach seinen Gitarrenkoffer aus dem VW-Bus holte.
Er setzte sich auf eine kurze Bank. Sie stand einladend unter einem
imposanten Lindenbaum, der den Garten gütig dominierte.

\sectionbreak

Da war sie also, seine neue Gitarre. Diesmal eine \emph{Gibson}. Wer
hätte das gedacht. Auch Meo hatte sich über seine Wahl gefreut. Der
Gedanke zauberte Louis ein Lächeln ins Gesicht. Sachte strich er über
die Saiten, klopfte den Korpus fein ab und ihm war, als halte er sein
Leben im Arm. Der Notizzettel in seiner Hosentasche beruhigte ihn.

Die Umgebung schien sich für Louis aufzulösen. Es wurde leer um ihn. Das
Treiben der andern empfand er als Film in Zeitlupe mit abgedrehtem Ton.
«L’amore è» würde warten können. Auf eine stille Zeit. Alleine
mit sich.

Louis begann das Instrument konzentriert zu stimmen. Die Klänge liessen
sein Herz augenblicklich hüpfen. Diese Freude hatte ihn lange alleine
gelassen und kehrte hier in dieser Nacht und in diesem Moment als kurze
Ahnung zu ihm zurück. Meos Familie war ihm ein reiner, praller
Glücksbringer.

\sectionbreak

Der Abend wurde zu einem musikalischen Highlight. Keiner hatte mehr
unangenehme Fragen gestellt. Die Zeit war rasend schnell vergangen. Was
immer sie spielten, sie fanden sich. Mit Leichtigkeit spielten sie alle
gleichzeitig, improvisierten, wechselten sich in kurzen und längeren
Soli ab, sangen einzeln und mehrstimmig. Es war ein Genuss.

Louis war sich sicher, sie hielten ihn endgültig für einen
professionellen Musiker. Spätestens nachdem er auf ihre Bitte, ein Solo
zu spielen, \emph{«Asturias»} von Albeniz dargeboten hatte, musste es
ihnen klar gewesen sein.

Meo und Vanessa war das Gitarrensolo bekannt.

\sectionbreak

Giorgia hatte es Louis eines Abends abgespielt und ausufernd geschwärmt.

«Ja, kenne ich, ist spitze, könnte ich auch,» hatte Louis ohne
nachzudenken gesagt. Sie hatte ihn erwartungsvoll angeschaut und
neckisch gelacht.

«Das will ich hören.»

Herausfordernd zwickte sie ihn in den Arm.

Es kostete ihn viel Übung, das Stück hinzubekommen. Er wollte ihr Held
sein.

\sectionbreak

Louis sass in ihrer Mitte. Sie legten ihre Instrumente zur Seite. Alle
Blicke richteten sich erwartungsvoll auf ihn. Der Glanz der Nacht hatte
sich über ihre Augen gelegt. Er spürte Nervosität aufkommen, blickte zu
Boden, schüttelte seine Hände und setzte an. Die ersten Töne erklangen.
Mit jedem Anschlag schwand seine Unsicherheit. Es gab nur noch ihn und
die Gitarre. Inmitten von friedvoller Zeitlosigkeit.

\qrblock{Isaac Albeniz, Simon Dinnigan \emph{«Asturias»} (Leyenda)}{media/QR-Codes/058_Asturias_(Leyenda)_Isaac_Albéniz,_Simon_Dinnigan_Spotify.png}

Louis schaute hoch. Die drei starrten ihn ungläubig an. Dann war, als
würden sie einen warmen Regen an Komplimenten über ihn giessen.

Die Worte von Meos Mutter verblüfften ihn besonders.

«Ich spiele ja kein Instrument, Ciro, darum kann ich nicht beurteilen,
was du da gerade gemacht hast,» ihre Augen glänzten, «aber dass das
unglaublich war, weiss ich. Grossartig, du hast meinen Respekt.»

Dann klatschten sie alle ein zweites Mal.

\sectionbreak

Später hatte Louis schliesslich nur noch gestaunt, als Meo sich an sein
vermeintlich gestohlenes Handy erinnerte. Und ihm sein altes,
«ausrangiertes», wie er sagte, schenkte. Er hatte Louis Handy und
Ladekabel in die Hände gelegt. Und nicht nur das.

«Ist frisch geladen, mein Freund,» hatte er mit Schalk hinzugefügt.

Louis hatte seine Hand auf Meos Schulter gelegt, als er sich bedankte.
Es fiel ihm schwer, Tränen der Rührung zurückzuhalten. Meo hatte diesen
aufrichtigen Blick, den Louis nicht oft gesehen hatte.

Dagegen waren all seine Lügen eine einzige Frechheit. Er war ein
schlechter Mensch. Auch dass er um Meos Handynummer bat, die er dann
sofort speichern und sich bei ihm melden würde, wie er Meo versicherte,
war ein reines Ablenkungstheater. Er würde es nicht tun.

Louis musste konsequent bleiben. Er musste seinen Plan durchziehen, den
Neubeginn mit Ciro schaffen. Morgen würde er als erstes eine Prepaid
Karte organisieren. Dann tauchte wieder Giorgias Bild vor ihm auf. Mein
Gott, Giorgia, was war mit ihr?

\sectionbreak

Es war noch dunkel, als Louis die Augen öffnete. Es hatte geklappt. Er
war vor Sonnenaufgang erwacht . Bereits am Abend zuvor hatte er sich mit
den Worten «und wenn ich am Morgen nicht mehr da bin, habe ich es
geschafft, frühzeitig aufzustehen» bei der Familie bedankt und sich
verabschiedet. Sie hatten gelacht. Witze gemacht. Offensichtlich
glaubten sie nicht, dass er so früh aufstehen würde.

Louis hatte alles für ein rasches Verschwinden bereit gelegt. Er
brauchte nur noch in die Kleider zu steigen.

\sectionbreak

Die nahegelegene Tankstelle in Urdorf war rund um die Uhr geöffnet. Es
mochte gegen halb sechs sein, schätzte Louis, als er das Geschäft
verliess. Die neue SIM-Karte in seiner Hand fühlte sich kostbar an.
Während er sich auf einer Bank neben der Tankstelle niederliess und den
ersten Bissen eines Sandwichs kaute, schaute er gebannt auf sein Handy.
Meo sei Dank. Er legte die Karte ein. Seine Hände zitterten vor
Aufregung. Das Tor zur Welt, irgendwie. Ihm wurde warm. Er richtete das
Gerät ein. Verblüfft starrte er auf die Uhrzeit auf dem Display. Zwanzig
vor fünf. Er war früher dran als gedacht. Erwartungsvoll lud er als
erstes eine Reihe Apps hinunter, auch Google Maps. Er zog den Zettel aus
seiner Hosentasche und gab die Hofadresse ein.

Da war es also. Der Hof lag leicht höher als \emph{Betten Dorf}. Er
würde zur Talstation \emph{Betten} reisen, und anschliessend auf die
\emph{Bettmeralp} aufsteigen. Sein Geld war knapp. Von Zürich musste es
bis \emph{Brig} reichen. Vielleicht lag sogar die Fahrt mit der
\emph{Matterhorn-Gotthard-Bahn} bis zur \emph{Talstation Betten} drin.
Er würde sich am Hauptbahnhof Zürich direkt informieren. Und Louis
versprach sich, Vanessa, Meo und seiner ganzen Familie seine Situation
aufrichtig zu schildern. Er würde sie besuchen. Zu gegebener Zeit.
Später.

Erst musste Maman von ihm hören. Und natürlich Marco.

Louis packte seine Sachen zusammen und ging los. Die Luft war feucht und
angenehm kühl. Die Strassen sahen erholt aus. Er war zuversichtlich.
Wenn es zügig ging, würde er den Bahnhof in gut zweieinhalb Stunden
erreichen. Es würde nach und nach heller werden. Auch wärmer. Es war
gut, rasch voranzukommen.

Aus unmittelbarer Nähe hörte er plötzlich eine aufgeregte, männliche
Stimme. Erst jetzt sah er die beiden. Einige Meter vor ihm auf dem
Gehsteig. Ein Mann stand gebückt neben einem andern, der am Boden lag.
Dieser lallte irgend etwas vor sich hin und fuchtelte mit den Armen.
Dann holte der Stehende aus und verpasste ihm einen Fusstritt.

Louis eilte auf sie zu und sein Ruf kam reflexartig.

«Hey.»

Der Stehende drehte sich torkelnd um. Hob die Faust und starrte ihn an.
Louis erschreckte sein verzerrtes Gesicht.

«Willst du, willst du auch eins auf die Fresse?»

Dann kam er näher auf ihn zu. Louis hob abwehrend die Hände.

«Nein, nein, kann ich dir helfen?»

Der Mann senkte ganz langsam seinen Arm. Seine erstaunten Augen schienen
ihm beinahe aus dem Gesicht zu fallen. Der Andere am Boden wimmerte vor
sich hin. Mit unbeholfenen Bewegungen wischte er sich Blut aus dem
Gesicht. Für einen Moment schien die Szene einzufrieren.

Das Motorengeräusch eines heranfahrenden Autos hinter Louis durchbrach
die Stille. Es kam näher. Louis blickte kurz hinter sich. Und schrak
zusammen. Ein Polizeiwagen. Nein. Er fühlte seine Beine vibrieren und
versteifte sich. Dann begann es. Seine Hände flatterten und gingen
unkontrolliert in ein Zittern über. Er presste seine Arme nahe an den
Körper.

\sectionbreak

\emph{Damals. Fribourg. Neun.}

\emph{Seit halb sechs liegt Louis wach an diesem Herbstmorgen. Es ist Sonntag
und sein Vater wird ihn heute abholen. Zu einer Männerfahrt, hat er
geheimnisvoll am Telefon gesagt, um neun. Nur ihn, ohne Lucia. Es wird
eine Überraschung, hat er hinzugefügt.}

\emph{Das Warten schmerzt Louis schon fast. Vielleicht darf er in einem
Oldtimer mitfahren. So wie vor ein paar Wochen, bevor Vater wieder einen
Wagen dem neuen Käufer übergab. Vater hat ganz viel Geld für diesen
orangen Alfa Romeo bekommen. Viele, viele Tausender, hat er Louis
erzählt.}

\emph{Louis frühstückt, sitzt herum und eilt alle fünf Minuten ans Fenster. Um
zehn nach neun ist der Vater noch nicht da. Louis ist ungemütlich. Das
fällt auch Maman auf. Er hüpft nervös herum, stört Lucia beim Malen und
weiss nicht mehr, wie er das Warten aushalten soll. Die Sekunden
schleichen dahin. Dann endlich ertönt ein Hupen von unten vor dem Haus.
Es ist sein Vater und halb zehn. Louis packt seine Jacke, küsst Maman,
ruft Lucia ein Tschüss zu und rennt aus der Tür. Atemlos unten
angekommen, reisst er die Eingangstür auf.}

\emph{Der Vater steht neben dem Auto. Er breitet seine Arme aus. Louis rennt
zu ihm und Vater hebt ihn hoch. Louis hat sich von Mamans Parfüm einen
Spritzer ins gegelte Haar gegeben. Er liebt es, wenn ihn sein Vater an
sich drückt und an ihm schnuppert. Vater verdreht dann die Augen so
lustig. «Ich fress dich auf, nimm dich in Acht,» sagt er dann, auch
jetzt, und kitzelt ihn. Und Louis kichert, bis ihm die Tränen kommen.}

\emph{«Nun los, steigen wir ein,» sagt Vater fröhlich. Louis staunt und
stutzt. «Der Ferrari, Papa, schön, aber…den habe ich doch schon
gesehen? Der war doch rot, nicht?» Sein Vater kneift sein linkes Auge
zu. «Ja, du Schlaumeier, aber der Käufer wollte einen schwarzen, darum,»
Vater weist mit einer weiten Armbewegung Richtung Fahrzeug, «darum habe
ich ihn umgespritzt, Tada, wie gefällt er dir?» Louis ist enttäuscht.
Rot hat er ihm besser gefallen. «Geht so,» sagt er, «rot war schöner.»
Aber eigentlich ist es egal. Louis freut sich, den Vater für sich
alleine zu haben.}

\emph{Louis sitzt neben Vater, natürlich vorne. Hinten ist eh kein Platz. Er
fühlt sich grossartig. Auch wenn er sich strecken muss, um aus der
Frontscheibe schauen zu können. Die Sitzerhöhung hat Vater vergessen. Es
riecht nach Leder wie auf Grandpères Sofa. Das Armaturenbrett glänzt.
Der Motor gibt ein röhrendes Brummen von sich. Das Geräusch ist
aufregend. Und als der Vater auf der Autobahn so richtig Gas gibt,
glaubt sich Louis im Paradies. Sie reden nicht viel. Louis geniesst. Die
Reise könnte bis ans Ende der Welt gehen. Er schaut seinen Vater immer
wieder verzückt von der Seite an. Er ist stolz auf ihn.}

\emph{Auf einmal verändert sich Vaters Gesicht. Louis sieht, dass er nervös
blinzelt, immer wieder in den Rückspiegel schaut und komisch atmet. «Was
ist los, Papa?» will Louis wissen, dreht sich um und sieht durchs
Heckfenster ein Polizeiauto mit einer beleuchteten Aufschrift. Dazu
schrillt das Martinshorn. Jetzt hat Louis Angst. «Papa, ein Polizeiauto,
da, da,» er will die verdrehten Buchstaben lesen, versucht es rückwärts,
«da steht «Stopp.» Louis hält sich die Ohren zu. Vater schreit. Er
klingt nun böse.}

\emph{«Sei still, Louis, und halt dich fest,» sagt er und drückt das Gaspedal
durch. Die Beschleunigung presst Louis ins Leder. Krampfhaft krallt er
seine Finger zu beiden Seiten in den Sitz. «Papa,» schreit er panisch.
Sie rasen auf der Überholspur an den Autos vorbei. Hinter ihnen folgt
der Polizeiwagen mit greller Sirene. Pure Angst hält Louis gefangen.
Vielleicht werden sie sterben, denkt er plötzlich. Er schliesst die
Augen und hält sie zusammengepresst. Er will nichts mehr sehen. Auf
einmal quietschen die Bremsen. Louis reisst die Augen auf, sieht das
Polizeiauto nun plötzlich vor ihnen, wird ruckartig nach vorne
geschlagen, hängt im Gurt und versucht sich am Armaturenbrett
abzustützen. Und alsbald steht ihr Auto still. Vater flucht und Louis
zittert. Sein Bauch tut weh. Ein Weinkrampf schüttelt seinen Körper. Ein
Polizist reisst die Führertür auf.}

\emph{«Kantonspolizei. Hände aufs Lenkrad,» ruft eine barsche Stimme in den
Wagen, «sind sie von Sinnen?»}

\emph{Der Polizist hat sich nun hinuntergebeugt und blickt direkt in Louis’
Augen, «und dann noch Fahrerflucht mit einem Kind. Steigen Sie aus,
sofort, bitte.» Der Polizist klingt bedrohlich wütend. «Da ist noch ein
Kind, Andrin, komm bitte!» ruft er nach vorn. Louis weint ununterbrochen
und muss geschehen lassen, was folgt. Die Wagen stehen nun
hintereinander auf dem Pannenstreifen. Der zweite Polizist hilft Louis
aus dem Auto und führt ihn zum Polizeiwagen. Er will wissen, ob der
Fahrer wirklich sein Vater ist. So etwas Dummes hat Louis noch keiner
gefragt.}

\emph{Der erste Polizist entfernt sich nach vorne an den Strassenrand. Nach
kurzer Zeit kommt er zurück.}

\emph{«Hier die Nummer der Mutter, Andrin, ruf sie bitte an. Sie kann den
Kleinen abholen,» sagt er, nachdem er mit dem Funkgerät in der Hand
wieder vor ihnen steht. Seinen Vater sieht Louis nicht mehr. Es ist
laut. Die Autos rasen links an ihnen vorüber. Der Polizist legt Louis
eine Decke über die Schulter und redet beruhigend auf ihn ein. Louis
hört etwas von «deine Mutter holt dich ab» und weint hemmungslos weiter.
Immer wieder ruft er nach seinem Vater. Aber der Polizist sagt, dass er
sich gedulden muss, und dass ihm nichts passieren wird und aufs Neue,
seine Mutter hole ihn bald ab. Louis ist froh, dass der Polizist bei ihm
bleibt. Er sitzt lange in die Decke gehüllt auf dem Rücksitz des
Polizeiwagens und wimmert vor sich hin.}

\emph{Irgendeinmal, nach ewig langem Warten, steht Maman vor der Autotür. Sie
hat sehr besorgte Augen. Und Louis schält sich aus der Decke, steigt aus
und wirft sich in ihre Arme. Sie hält ihn ganz lange eng an sich
gedrückt. Erst als er sich beruhigt, nimmt sie ihn an der Hand und sie
steigen in Mamans Auto. Louis macht sich grosse Sorgen um seinen Vater.
Das war eine richtige Verfolgungsjagd, wie er sie in Filmen gesehen hat.
Was wird nun mit Vater passieren? Louis’ Hände zittern immer noch. Er
kann nichts dagegen tun.}

\sectionbreak

Unzählige Male hatte es Louis geschafft. Hatte er sich Vater wieder ins
Lot träumen können. Dahin zurück, wo er sein perfekter Held war. Stark,
reich und allwissend.

Von diesem Tag an wurde es schwieriger.
